\documentclass[../DoAn.tex]{subfiles}
\begin{document}
\label{chapter:5}
Chương này trình bày các đóng góp trọng tâm và các giải pháp kỹ thuật cốt lõi mà đồ án đã thiết kế và triển khai nhằm giải quyết những thách thức đặc thù trong bài toán chatbot tư vấn pháp luật.
\section{Cơ chế kiểm tra tình trạng hiệu lực, xử lý phân loại các văn bản pháp luật (căn cứ chính, căn cứ bổ trợ, căn cứ hướng dẫn, sửa đổi, thay thế) và các căn cứ chồng chéo, xung đột}
\label{section:5.1}

\subsection{Dẫn dắt và giới thiệu bài toán}
\label{subsection:5.1.1}

Trong lĩnh vực tư vấn pháp luật Việt Nam, một thách thức lớn là tính thay đổi, cập nhật và thứ bậc của văn bản quy phạm pháp luật. Một văn bản luật có thể đúng tại thời điểm ban hành nhưng lại trở nên vô hiệu hoặc bị sửa đổi tại thời điểm người dùng đặt câu hỏi. Hơn nữa, một sự kiện pháp lý xảy ra trong quá khứ (ví dụ: một vụ việc ly hôn năm 2018) đòi hỏi phải áp dụng đúng văn bản có hiệu lực tại thời điểm đó, chứ không phải văn bản hiện hành. Ngoài ra, với cùng một vấn đề pháp lý, thực tế có thể tồn tại nhiều văn bản cùng điều chỉnh; khi đó cần áp dụng nguyên tắc hiệu lực và thứ bậc của các văn bản quy phạm pháp luật.

Vì vậy, đồ án tập trung vào cơ chế tự động kiểm tra hiệu lực theo mốc thời gian, truy vết phiên bản điều khoản trên đồ thị tri thức, đính kèm văn bản sửa đổi, hướng dẫn và tham chiếu, đồng thời cảnh báo khi ngữ cảnh truy xuất gồm các quy định đã bị thay thế hoặc hết hiệu lực. Cách tiếp cận này nhằm tăng độ tin cậy của câu trả lời và tạo bộ benchmark dataset để kiểm thử hệ thống với nhiều tình huống khác nhau.

\subsection{Giải pháp kỹ thuật: Truy xuất động trên Đồ thị tri thức theo thời gian}
\label{subsection:5.1.2}

Giải pháp được đề xuất dựa trên khả năng trích xuất ý định từ câu hỏi người dùng, lựa chọn căn cứ pháp lý và truy vấn đồ thị tri thức. Quy trình gồm ba giai đoạn chính: trích xuất và chuẩn hóa mốc thời gian; truy vết thực thể có hiệu lực trên đồ thị; chuẩn hóa ngữ cảnh phục vụ mô hình ngôn ngữ lớn (LLM).

\begin{table}[H]
\centering
\begin{tabular}{|p{15cm}|}
\hline
\textbf{Giải thuật:} Kiểm tra hiệu lực và chuẩn hóa ngữ cảnh pháp lý \\ \hline
\textbf{Đầu vào:} Câu hỏi người dùng $Q$ \\
\textbf{Đầu ra:} Ngữ cảnh pháp lý chuẩn hóa $C$ phục vụ sinh câu trả lời. \\ \hline

\textbf{Giai đoạn 1: Trích xuất và chuẩn hóa mốc thời gian} \\
1. $LLM\_Extract(Q) \rightarrow \{IDs, T_{event}\}$ \\
\textit{// Trích xuất danh sách căn cứ điều luật sẽ sử dụng từ tập điều luật định nghĩa sẵn theo dạng câu hỏi và mốc thời gian sự kiện (nếu có).} \\
2. Chuẩn hóa $T_{event}$ về dạng ngày đầy đủ (chỉ có năm $\rightarrow$ đầu năm; có năm-tháng $\rightarrow$ đầu tháng). \\
3. \textbf{Nếu} $T_{event} = \text{null}$: \\
\quad $T_{event} \leftarrow T_{today}$; $UserProvided \leftarrow \text{false}$. \\
\quad \textbf{Ngược lại}: $UserProvided \leftarrow \text{true}$. \\

\textbf{Giai đoạn 2: Truy vết thực thể có hiệu lực trên đồ thị tri thức} \\
4. Tìm các nút điều luật gốc có mã định danh thuộc $IDs$. \\
5. Mở rộng xuống các nút chi tiết (khoản, điểm) gắn với điều luật gốc. \\
6. Duyệt quan hệ thay thế theo cả hai chiều để thu toàn bộ phiên bản trên cùng dòng thay thế (quá khứ và tương lai). \\
7. Lọc phiên bản $V$ thỏa mãn hiệu lực tại $T_{event}$: \\
\quad $V.ngay\_co\_hieu\_luc \leq T_{event}$ và \\
\quad ($V.ngay\_het\_hieu\_luc = \text{null}$ hoặc $V.ngay\_het\_hieu\_luc > T_{event}$). \\
8. Với mỗi $V$, truy xuất văn bản sửa đổi, bổ sung $M$ qua quan hệ DUOC\_SUA\_DOI\_BOI, chỉ giữ $M$ còn hiệu lực tại $T_{event}$. \\
9. Truy xuất văn bản hướng dẫn $G$ qua quan hệ HUONG\_DAN\_BOI; mở rộng chi tiết hướng dẫn; kiểm tra sửa đổi đối với chi tiết hướng dẫn (nếu có). \\
10. Truy vết theo chiều thay thế về phía tương lai để xác định quy định hiện hành $V_{current}$ phục vụ đối chiếu. \\
11. Truy xuất các quy định tham chiếu bổ trợ liên quan, có lọc hiệu lực tại $T_{event}$. \\
12. Gói kết quả thô thành tập dữ liệu $D = \{\text{căn cứ chính}, \text{hướng dẫn}, \text{tham chiếu}, \text{quy định đối chiếu}\}$. \\

\textbf{Giai đoạn 3: Chuẩn hóa và gán nhãn ngữ cảnh cho LLM} \\
13. Duyệt căn cứ chính: nếu tồn tại $M$, thay nội dung hiển thị bằng nội dung sửa đổi nhưng giữ nhãn văn bản gốc; kích hoạt cờ \texttt{SUA\_DOI}. \\
14. Phân loại căn cứ hướng dẫn và căn cứ tham chiếu bổ trợ; xử lý sửa đổi đối với hướng dẫn tương tự căn cứ chính. \\
15. Kích hoạt cờ \texttt{UU\_TIEN} nếu trong $D$ xuất hiện nhiều hơn một giá trị cấp bậc pháp lý. \\
16. Kích hoạt cờ \texttt{LICH\_SU} khi $UserProvided = \text{true}$, căn cứ tại $T_{event}$ đã hết hiệu lực và tồn tại $V_{current}$. \\
17. Tổng hợp thông tin hiệu lực theo tên văn bản (ngày có hiệu lực, ngày hết hiệu lực). \\
18. Sắp xếp căn cứ chính và hướng dẫn theo cấp bậc tăng dần (số càng nhỏ, cấp càng cao). \\
19. Đóng gói $C$ thành chuỗi có cấu trúc: thông tin cảnh báo (các cờ cảnh báo để LLM biết cách trả lời), hiệu lực, căn cứ chính, căn cứ tham chiếu bổ trợ, căn cứ hướng dẫn và văn bản thay thế. \\ \hline
\end{tabular}
\caption{Mô tả giải thuật truy xuất và phân loại các văn bản pháp lý}
\label{table:algo_temporal}
\end{table}

Dựa trên Bảng~\ref{table:algo_temporal}, ở Giai đoạn~1 hệ thống gán các biến \texttt{target\_ids} và \texttt{target\_date} từ kết quả phân tích của tác nhân. Mốc thời gian được chuẩn hóa trước khi truy vấn; nếu người dùng không nêu thời điểm sự kiện, hệ thống tự động dùng ngày hiện tại (\texttt{today\_date}) làm căn cứ truy xuất và đánh dấu người dùng chưa cung cấp mốc thời gian cụ thể.

Giai đoạn~2 là đóng góp quan trọng nhất về mặt giải thuật. Hệ thống không chỉ truy xuất nút điều luật ban đầu mà còn mở rộng xuống các nút chi tiết (khoản, điểm), sau đó duyệt toàn bộ dòng thay thế trên đồ thị theo cả hai chiều thời gian. Điều kiện lọc hiệu lực được áp dụng thống nhất cho phiên bản áp dụng, văn bản sửa đổi, hướng dẫn và tham chiếu. Song song, thuật toán truy vết quy định hiện hành trên cùng dòng thay thế để phục vụ đối chiếu khi người dùng hỏi theo mốc thời gian quá khứ. Nhờ đó, hệ thống đạt được: (i)~nhận diện phiên bản điều khoản có hiệu lực tại thời điểm sự kiện; (ii)~liên kết văn bản sửa đổi, bổ sung và hướng dẫn trực tiếp với phiên bản đang xét; (iii)~bổ sung quy định hiện hành làm căn cứ đối chiếu. Toàn bộ căn cứ thỏa điều kiện hiệu lực và liên kết đã khai báo trên đồ thị được đưa vào tập dữ liệu thô để xử lý ở bước sau.

Ở Giai đoạn~3, module chuẩn hóa ngữ cảnh phân loại dữ liệu thô thành: (i)~\textbf{căn cứ chính}: các điều khoản trực tiếp áp dụng; (ii)~\textbf{căn cứ tham chiếu bổ trợ}: quy định được tham chiếu từ căn cứ chính, hoặc được tham chiếu từ các căn cứ hướng dẫn hoặc sửa đổi; (iii)~\textbf{căn cứ hướng dẫn}: văn bản hướng dẫn thi hành; (iv)~\textbf{văn bản thay thế}: các điều luật hiện hành thay thế; (v)~\textbf{thông tin cảnh báo}: các cờ \texttt{LICH\_SU}, \texttt{UU\_TIEN} và \texttt{SUA\_DOI} cùng khối thông tin hiệu lực văn bản. Cờ \texttt{UU\_TIEN} chỉ kích hoạt khi ngữ cảnh chứa nhiều cấp bậc pháp lý khác nhau; đây là cảnh báo thứ bậc nhắc LLM cân nhắc nguyên tắc ưu tiên khi trình bày. Cờ \texttt{LICH\_SU} được kích hoạt khi người dùng nêu mốc thời gian quá khứ, căn cứ truy xuất đã hết hiệu lực tại mốc đó và tồn tại quy định thay thế hiện hành, giúp LLM thông báo cho người dùng về sự thay đổi pháp luật.

Về chuẩn hóa nội dung cho LLM: thay vì đồng thời đưa cả nội dung gốc và nội dung sửa đổi gây nhiễu, thuật toán ưu tiên hiển thị nội dung từ văn bản sửa đổi nhưng vẫn giữ tên căn cứ gốc, theo dạng ``Theo quy định tại [văn bản gốc], được sửa đổi, bổ sung bởi [văn bản mới]\ldots''. Danh sách căn cứ chính và hướng dẫn được sắp xếp theo cấp bậc tăng dần (Luật trước Nghị định, Thông tư) nhằm hướng dẫn LLM trích dẫn các văn bản pháp lý theo đúng thứ tự. Cuối cùng, ngữ cảnh được đóng gói thành các mục có tiêu đề rõ ràng (bao gồm thông tin cảnh báo, thông tin hiệu lực các văn bản, căn cứ chính, căn cứ tham chiếu bổ trợ, căn cứ hướng dẫn và văn bản thay thế) trước khi đưa vào prompt sinh câu trả lời.

\subsection{Xây dựng tập Benchmark đánh giá hệ thống tư vấn pháp luật}
\label{subsection:benchmark}

Ngoài cơ chế truy xuất theo thời gian trên đồ thị tri thức, đồ án đóng góp một tập benchmark dataset phục vụ đánh giá có hệ thống chất lượng câu trả lời của chatbot trong lĩnh vực Luật Hôn nhân và gia đình. Tập dữ liệu được xây dựng qua các giai đoạn như sau:

\subsubsection{Thu thập dữ liệu}

Nguồn dữ liệu ban đầu lấy từ chuyên mục \textit{Hỏi đáp pháp luật} trên cổng thuvienphapluat.vn, được tổ chức theo dạng câu hỏi pháp lý, ví dụ: điều kiện kết hôn, đăng ký kết hôn, chế độ tài sản vợ chồng, ly hôn, cấp dưỡng\ldots

Trong quá trình crawl dữ liệu, khó khăn phát sinh vì các câu hỏi, câu trả lời và đặc biệt là phần căn cứ pháp lý của mỗi câu hỏi trên trang thuvienphapluat.vn được trình bày theo ba kiểu HTML khác nhau tùy theo từng bài đăng. Script crawl dữ liệu phải xử lý và nhận diện được cả ba trường hợp này để lấy về đúng những thông tin cần thiết. Các thông tin được lấy về ở giai đoạn crawl gồm question (câu hỏi của người dùng) và ground\_truth (câu trả lời nguyên văn từ đội tư vấn Thư viện Pháp luật). Sau khi thu thập các câu hỏi đã được phân loại theo từng chủ đề ở chuyên mục hỏi đáp pháp luật, cần tiếp tục phân loại lại dạng (chủ đề) câu hỏi theo cách thủ công. Nguyên nhân là tuy mục hỏi đáp pháp luật của thuvienphapluat.vn đã được chia thành các chủ đề, các bài đăng thuộc cùng một chủ đề vẫn có thể chứa câu hỏi thuộc chủ đề khác. Ví dụ, bài đăng thuộc chủ đề điều kiện kết hôn vẫn có thể chứa câu hỏi về kết hôn trái pháp luật và ngược lại. Vì muốn chia và đánh giá các câu hỏi trong tập benchmark dataset theo chủ đề, đồ án phải thực hiện bước phân loại lại. Ở bước này có thể sử dụng LLM để phân loại câu hỏi theo chủ đề, nhưng khi xử lý số lượng lớn câu hỏi được lấy về (khoảng 700 câu), phương pháp này phát sinh chi phí lớn và vẫn cần đánh giá thủ công lại.

Quá trình thu thập các câu hỏi và câu trả lời mẫu từ các trang web khác không gặp quá nhiều khó khăn vì định dạng câu hỏi, câu trả lời có bố cục rõ ràng.

\subsubsection{Cấu trúc mẫu dữ liệu trong tập benchmark}

Mỗi mẫu trong tập benchmark dataset là một đối tượng JSON gồm 11 trường chuẩn. Bảng~\ref{tab:benchmark_schema} mô tả tên trường, kiểu dữ liệu và ý nghĩa của từng trường trong một mẫu benchmark.

\begin{table}[H]
\centering
\caption{Cấu trúc trường dữ liệu của một mẫu benchmark}
\label{tab:benchmark_schema}
\begin{tabular}{|p{4.5cm}|p{2.5cm}|p{7.5cm}|}
\hline
\textbf{Tên trường} & \textbf{Kiểu dữ liệu} & \textbf{Ý nghĩa} \\
\hline
question & String & Câu hỏi của người dùng (thu từ nguồn crawl hoặc biên soạn bổ sung). \\
\hline
ground\_\allowbreak truth & String & Câu trả lời tham chiếu từ đội ngũ tư vấn pháp luật Thư viện Pháp luật; dùng làm chuẩn so sánh khi đánh giá chatbot. \\
\hline
chu\_\allowbreak de\_\allowbreak phap\_\allowbreak ly & String & Mã chủ đề pháp lý tương ứng retriever/agent (vd.\ che\_do\_tai\_san\_cua\_vo\_chong, ly\_hon\ldots). \\
\hline
loai\_\allowbreak cau\_\allowbreak hoi & String & Phân loại dạng câu hỏi, thường là \textit{Lý thuyết} hoặc \textit{Tình huống}. \\
\hline
van\_\allowbreak ban\_\allowbreak phap\_\allowbreak luat\_\allowbreak lien\_\allowbreak quan & List & Tập mã định danh văn bản pháp luật có liên quan đến câu hỏi (vd.\ Luat\_HNGD\_2014, BoLuat\_DanSu\_2015). \\
\hline
tinh\_\allowbreak trang\_\allowbreak hieu\_\allowbreak luc & Dictionary & Ánh xạ từng điều/khoản (khóa ID) sang trạng thái hiệu lực tại thời điểm đánh giá (vd.\ CON\_HIEU\_LUC). \\
\hline
can\_\allowbreak cu\_\allowbreak phap\_\allowbreak ly\_\allowbreak chinh & List & Danh sách ID điều/khoản là căn cứ pháp lý chính bắt buộc phải được truy xuất và trích dẫn. \\
\hline
can\_\allowbreak cu\_\allowbreak phap\_\allowbreak ly\_\allowbreak bo\_\allowbreak tro & List & Căn cứ tham chiếu bổ trợ (điều luật/văn bản được dẫn chiếu thêm, không phải trọng tâm). \\
\hline
van\_\allowbreak ban\_\allowbreak huong\_\allowbreak dan\_\allowbreak sua\_\allowbreak doi\_\allowbreak bo\_\allowbreak sung\_\allowbreak thay\_\allowbreak the & List & Văn bản hướng dẫn, sửa đổi, bổ sung hoặc có liên quan thay thế cần đính kèm khi trả lời. \\
\hline
can\_\allowbreak cu\_\allowbreak khong\_\allowbreak nen\_\allowbreak ap\_\allowbreak dung & List & Các căn cứ không được áp dụng tại thời điểm sự kiện pháp lý: thường là văn bản thuộc quá khứ (đã hết hiệu lực/bị thay thế) hoặc tương lai (chưa có hiệu lực) so với mốc thời gian của tình huống. \\
\hline
quan\_\allowbreak he\_\allowbreak chong\_\allowbreak cheo\_\allowbreak mau\_\allowbreak thuan & List & Ghi nhận quan hệ chồng chéo hoặc mâu thuẫn giữa các căn cứ (nếu có), phục vụ đánh giá khả năng phát hiện xung đột ngữ cảnh. \\
\hline
\end{tabular}
\end{table}

\subsubsection{Quy trình hiệu chỉnh và gán nhãn thủ công}

Sau khi crawl, tập benchmark trải qua các bước hiệu chỉnh sau:

\textbf{Bước 1: Gán chủ đề pháp lý}
Gán chu\_de\_phap\_ly cho từng câu hỏi; loại bỏ hoặc chỉnh sửa mẫu nhiễu (nội dung không thuộc HNGĐ, trùng lặp, lỗi phân tích cú pháp).

\textbf{Bước 2 --- Cập nhật căn cứ pháp lý cho câu hỏi không có dữ kiện thời gian.}
Đối với các câu hỏi không nêu mốc thời gian sự kiện, hệ thống chatbot mặc định truy xuất theo thời điểm hiện tại. Do đó, trong tập benchmark, người hiệu chỉnh phải tự cập nhật các căn cứ pháp lý trong ground\_truth và các trường metadata liên quan (can\_cu\_phap\_ly\_chinh, van\_ban\_phap\_luat\_lien\_quan, tinh\_trang\_hieu\_luc\ldots) theo phiên bản văn bản mới nhất còn hiệu lực, thay cho các viện dẫn luật cũ có thể còn sót trong bài tư vấn gốc.

\textbf{Bước 3 --- Gán nhãn pháp lý dựa trên ground truth.}
Dựa trên nội dung ground\_truth từ đội ngũ tư vấn pháp luật, người hiệu chỉnh điền thủ công các trường:
\begin{itemize}
    \item loai\_cau\_hoi: phân biệt câu lý thuyết và câu tình huống cụ thể;
    \item van\_ban\_phap\_luat\_lien\_quan, tinh\_trang\_hieu\_luc: ID của các văn bản liên quan và trạng thái hiệu lực từng điều khoản;
    \item can\_cu\_phap\_ly\_chinh, can\_cu\_phap\_ly\_bo\_tro, \\ van\_ban\_huong\_dan\_sua\_doi\_bo\_sung\_thay\_the: tách căn cứ chính, bổ trợ, hướng dẫn;
    \item can\_cu\_khong\_nen\_ap\_dung: liệt kê các ID văn bản không phù hợp với thời điểm xảy ra sự kiện---bao gồm văn bản quá khứ đã bị thay thế/hết hiệu lực hoặc văn bản tương lai chưa áp dụng---nhằm đánh giá khả năng hệ thống không đưa nhầm căn cứ lỗi thời vào câu trả lời;
    \item quan\_he\_chong\_cheo\_mau\_thuan: ghi nhận xung đột/chồng chéo giữa căn cứ (nếu phát hiện khi đối chiếu).
\end{itemize}

\subsubsection{Thiết kế test case cho các trường hợp ngoại lệ}

Bên cạnh các mẫu Q\&A chuẩn từ nguồn tư vấn, đồ án dự kiến bổ sung nhóm test case ngoại lệ để kiểm tra và đánh giá quá trình tool selection của router agent, quá trình retrieve của retriever agent và các test case có các câu hỏi thuộc trường hợp đặc biệt (ví dụ chào hỏi, hỏi mình có thể làm gì, hỏi câu không liên quan luật hôn nhân gia đình. Những câu hỏi có liên quan nhưng không có thông tin, ví dụ hỏi tình huống năm 1998. Test các câu hỏi đòi hỏi aggregations và filtering (Có bao nhiêu điều? Có bao nhiêu nghị định ...) Cụ thể như sau:

\subsection{Kết quả đạt được}
\label{subsection:5.1.3}

Giải pháp kiểm tra hiệu lực động đã mang lại những cải thiện về chất lượng tư vấn: (i)Tính minh bạch cao: hệ thống không chỉ trả lời kết quả mà còn giải thích được ``tại sao văn bản này được chọn'' dựa trên mốc thời gian; (ii)khắc phục hiện tượng ảo giác (Hallucination): bằng cách giới hạn LLM chỉ được suy luận trên ``ngữ cảnh sạch'' đã qua lọc hiệu lực từ Neo4j, tỷ lệ viện dẫn sai văn bản đã hết hiệu lực giảm xuống dưới 10\% (theo chỉ số Invalid Citation Rate); (iii)xử lý xung đột pháp lý: hệ thống đạt độ chính xác trên 80\% trong việc nhận diện và cảnh báo các trường hợp văn bản chồng chéo. Người dùng nhận được thông tin đầy đủ về việc một điều luật được sửa đổi bởi nghị định nào, giúp tăng tính thuyết phục của phản hồi.

\end{document}